% SOURCE_STATEMENT: "For $i=1,2$, let $$\\Lambda_i[X]=\\sum_{\\mu=0}^{d_i^2-1}\\lambda_\\mu^{(i)}F_\\mu^{(i)}XF_\\mu^{(i)},$$ where the $F_\\mu^{(i)}$ form Hermitian, Hilbert–Schmidt orthonormal bases of $M_{d_i}(\\mathbb{C})$. If exactly one coefficient of $\\Lambda_2$ is negative and every positive coefficient is at least its absolute value, then $\\Lambda=\\Lambda_1\\otimes\\operatorname{id}_{d_2}+\\operatorname{id}_{d_1}\\otimes\\Lambda_2$ is positive. Can this family be generalized by relaxing the orthonormal-basis hypothesis? For $$\\rho_I=\\frac1{|I|}\\sum_{(\\alpha,\\beta)\\in I}|\\psi_{\\alpha\\beta}\\rangle\\langle\\psi_{\\alpha\\beta}|,$$ determine the properties of these states and, in particular, whether their entanglement can be detected by other criteria."
% Preparation and final self-audit: CONTRIBUTING.md prompts C and D; source/audit artifacts retained in the local physics_statements_2026_10_11 batch.
% !TeX program = lualatex
\documentclass[11pt]{article}
\usepackage[margin=25mm]{geometry}
\usepackage{fontspec}
\setmainfont{Latin Modern Roman}
\usepackage{amsmath,amssymb,mathtools}
\usepackage{unicode-math}
\setmathfont{Latin Modern Math}
\usepackage{xurl,hyperref}
\hypersetup{colorlinks=true,urlcolor=blue}
\setcounter{secnumdepth}{0}
\setlength{\parindent}{0pt}
\setlength{\parskip}{5pt}
\setlength{\emergencystretch}{3em}
\newcommand{\sref}[2]{\hyperlink{source:#1:#2}{[S#2]}}
\newcommand{\cataloguescope}[1]{\paragraph{Recorded target.}#1}
\newcommand{\R}{\mathbb R}
\newcommand{\C}{\mathbb C}
\newcommand{\N}{\mathbb N}
\newcommand{\Z}{\mathbb Z}
\newcommand{\Q}{\mathbb Q}
\newcommand{\F}{\mathbb F}
\newcommand{\eps}{\varepsilon}
\begin{document}
\section{Positive Maps and Multipartite Bound Entanglement}\label{statement}

\subsection{Definitions and mathematical statement}

\paragraph{Shared notation.}$\mathbb N=\{1,2,\ldots\}$; the standard number fields have their usual meanings. All additional conventions are stated locally.


\paragraph{Definitions and conventions.} $M_d(\mathbb C)$ is the matrix algebra, with Hilbert--Schmidt inner product $\langle F,G\rangle=\operatorname{Tr}(F^\dagger G)$. A map is positive if it sends positive semidefinite matrices to positive semidefinite matrices, and completely positive if every tensor amplification by an identity map is positive. The source uses Hermitian orthonormal bases $F^{(i)}_\mu$ and real coefficients $\lambda^{(i)}_\mu$. For the four-qubit lattice states, let $\sigma_0=I_2$ and $\sigma_1,\sigma_2,\sigma_3$ be Pauli matrices; set $|\psi_+\rangle=\tfrac12\sum_{a,b=1}^2|a,b\rangle_A|a,b\rangle_B$ and $|\psi_{\alpha\beta}\rangle=(I_4\otimes\sigma_\alpha\otimes\sigma_\beta)|\psi_+\rangle$. Here $\varnothing\ne I\subseteq\{0,1,2,3\}^2$ and $|I|$ is its cardinality. PPT means positivity after partial transposition across the $A|B$ partition.

\paragraph{Source scope and qualifications.} The supplied positivity preamble omits the coefficient hypotheses on Lambda\_1 that are needed for the source's proposition. An arbitrary Lambda\_1 with large negative coefficients does not satisfy the stated conclusion. The primary source requires all coefficients of Lambda\_1 to be positive and sufficiently large relative to the single negative coefficient of Lambda\_2. This missing premise is explicitly restored in the source clarification.

\paragraph{Primary-source hypothesis.} In the report's positivity proposition, all $\lambda^{(1)}_\mu$ are positive and are at least the absolute value of the single negative coefficient of $\Lambda_2$; every other coefficient of $\Lambda_2$ has the same lower bound. The request is to generalize this family, for example by relaxing the orthonormal-basis identity, and to study the entanglement-detection properties of the states $\rho_I$.

\paragraph{Catalogue mathematical target.}
\paragraph{Catalogue wording retained for traceability.}
For $i=1,2$, let $$\Lambda_i[X]=\sum_{\mu=0}^{d_i^2-1}\lambda_\mu^{(i)}F_\mu^{(i)}XF_\mu^{(i)},$$ where the $F_\mu^{(i)}$ form Hermitian, Hilbert–Schmidt orthonormal bases of $M_{d_i}(\mathbb{C})$. If exactly one coefficient of $\Lambda_2$ is negative and every positive coefficient is at least its absolute value, then $\Lambda=\Lambda_1\otimes\operatorname{id}_{d_2}+\operatorname{id}_{d_1}\otimes\Lambda_2$ is positive. Can this family be generalized by relaxing the orthonormal-basis hypothesis? For $$\rho_I=\frac1{|I|}\sum_{(\alpha,\beta)\in I}|\psi_{\alpha\beta}\rangle\langle\psi_{\alpha\beta}|,$$ determine the properties of these states and, in particular, whether their entanglement can be detected by other criteria.

\paragraph{Saved literature context.} The supplied catalogue's assessment dated 2026-08-21 records: The report identifies open generalizations of the positive-map family and questions about the PPT bound-entangled states rho-tilde-I. The checked literature gives important criteria and constructions but no complete resolution of these exact questions. The assessment is retained as dated provenance; it is not a new comprehensive literature or open-status certification.
\paragraph{Scope and source qualification.}The supplied positivity preamble omits the coefficient hypotheses on Lambda_1 that are needed for the source's proposition. An arbitrary Lambda_1 with large negative coefficients does not satisfy the stated conclusion. The primary source requires all coefficients of Lambda_1 to be positive and sufficiently large relative to the single negative coefficient of Lambda_2. This missing premise is explicitly restored in the source clarification.
\subsection{Short English statement}
Extend the report's family of positive maps and determine how the associated Bell-diagonal states can exhibit detectable bound entanglement.
\subsection{Sources}
\begin{itemize}
\item[\hypertarget{source:statement:1}{S1}] ULAM.AI and the UnsolvedMath Contributors, UnsolvedMath, catalogue entry. Catalogue curation under CC BY 4.0; definitions and English exposition added locally.
\url{https://huggingface.co/datasets/ulamai/UnsolvedMath}.
\item[\hypertarget{source:statement:2}{S2}] Entanglement and Decoherence: Mathematics and Physics of Quantum Information and Computation (2005), Oberwolfach Reports, DOI 10.4171/owr/2005/04, pp. 185--255; official report source and mathematical context.
\url{https://doi.org/10.4171/owr/2005/04}.
\item[\hypertarget{source:statement:3}{S3}] Official EMS Press PDF of the same Oberwolfach report; PDF pages 48--50, printed pages 232--234.
\url{https://ems.press/content/serial-article-files/45978}.
\item[\hypertarget{source:statement:4}{S4}] Horodecki et al., multipartite separability and positive maps; reference in the saved dated catalogue literature review.
\url{https://arxiv.org/abs/quant-ph/0006071}.
\end{itemize}
\label{end:statement}
\end{document}
